\section{Limitations and Conclusion}

This work audited whether a frozen visual anomaly screen returns multiscale
evidence to the base-grid cells that generated it. The released ViT4TS
projector shifts every valid pooled-scale query by one flattened cell and
leaves the terminal query unsupported. IHP provides the minimal repair: it
uses the supplied zero-based memberships literally and then applies the
unchanged harmonic projector on the corrected support. In the 492-series
same-cache comparison, IHP improves equal-subdataset F1-max, AUPRC, and VUS-PR
by 0.0267, 0.0218, and 0.0102. Unadjusted intervals exclude zero for AUPRC
and VUS-PR but not F1-max. The paper-reported ViT4TS value of 0.612 and IHP's
0.662 provide secondary external context, not a paired method effect.

The evidence is bounded to the univariate ViT4TS screening stage, its released
$14\times14$ grid and pooling masks, and three benchmark families. IHP was a
prespecified component promoted after a composite candidate failed on the same
evaluation, so its intervals are not selection-adjusted. Moreover, released
full-series preprocessing and all-window median memory make the evaluated
screen offline and transductive; IHP does not remove this deployment
limitation. We do not evaluate downstream VLM verification, multivariate
renderers, online drift, or other visual backbones. Future work should make
coordinate contracts executable across other renderers and test them on an
independent, streaming-compatible protocol. Within the present boundary, the
exact 196/196 coverage certificate and paired ranking gains show that auditing
an inverse-coordinate interface can improve a frozen screen without training,
labels, learned parameters, or another encoder pass.
